%---------------------------------------------------------------------
% Urdu edition: title page, how to use, exam map.
%---------------------------------------------------------------------
\thispagestyle{empty}

% The band has to clear the whole masthead. Nastaliq at Scale=1.45 with
% \linespread{1.55} sets far taller lines than the nominal point sizes
% suggest, so this depth is measured from the rendered page, not computed.
\begin{tikzpicture}[remember picture, overlay]
  \fill[bmblued] (current page.north west) rectangle
        ([yshift=-12.6cm]current page.north east);
  \fill[bmaccent] ([yshift=-12.6cm]current page.north west) rectangle
        ([yshift=-12.82cm]current page.north east);
\end{tikzpicture}

\vspace*{1.1cm}
{\color{white}
  {\large علامہ اقبال اوپن یونیورسٹی، اسلام آباد}\par
  \vspace{4pt}
  {شعبۂ ریاضی \ \textbullet\ \ کلیۂ سائنس}\par
  \vspace{22pt}
  {\Huge\bfseries کاروباری ریاضی}\par
  \vspace{10pt}
  {\LARGE\bfseries متوقع سوالات اور ان کے حل}\par
  \vspace{16pt}
  {\large کورس کوڈ \textenglish{1429} \ (بعض اوقات \textenglish{5405}) \ \textbullet\ \ یونٹ \textenglish{1} تا \textenglish{9}}\par
  \vspace{6pt}
  % The \par must come BEFORE the group closes: colour is applied when the
  % paragraph is broken, so ending the group first reverts to black ink and
  % the line disappears into the dark band.
  {بی اے \ \textbullet\ \ بی کام \ \textbullet\ \ اے ڈی سی \ \textbullet\ \ بی ایس\par}
}

\vspace{1.4cm}

\begin{tcolorbox}[enhanced, colback=bmtint, colframe=bmblue, boxrule=0.6pt,
                  arc=3pt, left=14pt, right=14pt, top=12pt, bottom=12pt]
\textbf{\color{bmblue}یہ کتابچہ کیا ہے؟}\ \
یہ اے آئی او یو کی سرکاری درسی کتاب \emph{کاروباری ریاضی} (کورس کوڈ
\textenglish{5405/1429}، شعبۂ ریاضی، اے آئی او یو اسلام آباد، اشاعت
\textenglish{2023}) کے ہر یونٹ کے آخر میں دیے گئے
\textenglish{Self Assessment Questions} میں سے سب سے اہم اور متوقع سوالات کا
مختصر حل ہے۔

\medskip
\textbf{\color{bmblue}یہی سوالات کیوں؟}\ \
یونیورسٹی اپنے سمسٹر کے اسائنمنٹ اور فائنل پرچہ اسی ذخیرے سے بناتی ہے۔ اس لیے
کتاب کے انہی سوالات کی مشق دراصل پرچے ہی کی مشق ہے۔

\medskip
\textbf{\color{bmblue}تفصیلی ایڈیشن۔}\ \
تمام نو یونٹوں کے مکمل حل (انگریزی، \textenglish{84} صفحات) الگ فائل
\textenglish{AIOU\_1429\_Solved\_QuestionBank.pdf} میں موجود ہیں۔ یہ اردو
کتابچہ اسی کا مختصر خلاصہ ہے۔

\medskip
\textbf{\color{bmblue}حیثیت۔}\ \
یہ ذاتی مطالعے کے لیے تیار کیا گیا ہے۔ یہ اے آئی او یو کی مطبوعہ کتاب نہیں اور
نہ ہی اسے یونیورسٹی کی توثیق حاصل ہے۔ اسے من و عن نقل کر کے اسائنمنٹ میں جمع
کرانا منع ہے۔
\end{tcolorbox}

\clearpage

%---------------------------------------------------------------------
\section*{اسے کیسے استعمال کریں}
\markboth{اسے کیسے استعمال کریں}{}

سوال پڑھیے۔ حل کو ہاتھ یا کاغذ سے ڈھانپ لیجیے۔ پہلے خود کاپی پر حل کیجیے۔
\emph{اس کے بعد} ملائیے۔ صرف حل پڑھ لینے سے سیکھنے کا گمان تو ہوتا ہے مگر
امتحان میں کام وہی آتا ہے جو ناکام کوشش کے بعد سمجھ آیا ہو۔

\medskip
اس کتابچے میں تین چیزیں اسی مقصد کے لیے رکھی گئی ہیں:

\begin{itemize}
  \item \textbf{جواب کا خانہ (سبز)} --- آخری نتیجہ الگ نمایاں ہے، تاکہ آپ
        طریقہ پڑھے بغیر صرف اپنا جواب ملا سکیں۔
  \item \textbf{پڑتال کی سطر} --- جواب کو اصل مساوات میں واپس رکھ کر جانچا گیا
        ہے۔ اے آئی او یو طریقۂ کار کے نمبر دیتی ہے، اور پڑتال سے علامت
        (\textenglish{sign}) کی غلطی ممتحن سے پہلے آپ کو پکڑ میں آ جاتی ہے۔
  \item \textbf{عام غلطی (سرخ)} --- اس قسم کے سوال میں بار بار ہونے والی مخصوص
        غلطی۔ یہ خانے حل سے زیادہ قیمتی ہیں۔
\end{itemize}

\begin{ghalti}[اسائنمنٹ میں نقل نہ کریں]
اے آئی او یو سرقہ (\textenglish{plagiarism}) کی جانچ کرتی ہے، اور یہ حل عام
دستیاب ہیں۔ ان کا مقصد صرف یہ ہے کہ آپ اپنا کام جانچ سکیں اور وہ قدم پہچان سکیں
جہاں آپ کا طریقہ ٹوٹتا ہے۔ طریقہ سیکھیے، حساب اپنا لکھیے۔
\end{ghalti}

%---------------------------------------------------------------------
\section*{پرچے میں کس یونٹ کا کتنا وزن ہے}

پرچہ نو کے نو یونٹوں سے آتا ہے، مگر برابر نہیں۔ ذیل کی ترتیب کتاب کے اپنے وزن
اور گردش کرتے پرانے پرچوں سے مطابقت رکھتی ہے۔

\medskip
\begin{center}
\begin{tabular}{@{}p{4.6cm} p{2.4cm} p{6.4cm}@{}}
\toprule
\textbf{یونٹ اور موضوع} & \textbf{وزن} & \textbf{اصل امتحان کس چیز کا ہے} \\
\midrule
\textenglish{1} --- احتمال            & زیادہ   & جمع کا اصول بمقابلہ ضرب کا اصول \\
\textenglish{2} --- بے ترتیب متغیر    & درمیانہ & احتمالات کا مجموعہ ہمیشہ \textenglish{1} \\
\textenglish{3} --- مساوات و عدم مساوات & درمیانہ & علامت کا الٹنا؛ مطلق قدر کی دو صورتیں \\
\textenglish{4} --- خطی مساوات        & زیادہ   & پہلے میلان، پھر ایک نقطہ \\
\textenglish{5} --- میٹرکس            & زیادہ   & ضرب کے لیے ابعاد کا میل \\
\textenglish{6} --- محدد اور معکوس    & زیادہ   & ہم عاملوں کی علامت؛ کریمر کا اصول \\
\textenglish{7} --- تفرق              & زیادہ   & تقسیم کے اندر زنجیری اصول \\
\textenglish{8} --- جزوی مشتق         & درمیانہ & دوسرے متغیر کو ثابت رکھنا \\
\textenglish{9} --- بہتر بندی         & زیادہ   & دوسرے مشتق کی جانچ، ہمیشہ \\
\bottomrule
\end{tabular}
\end{center}

\medskip
\noindent
یونٹ \textenglish{4}، \textenglish{5}، \textenglish{6}، \textenglish{7} اور
\textenglish{9} مل کر زیادہ تر نمبر اور تمام طویل سوالات اٹھاتے ہیں۔ اگر وقت کم
ہو تو تیاری اسی ترتیب سے کیجیے۔

\clearpage
